% Глава: 18-ipc
% Статус: наполнено (первый прогон).
\chapter{Взаимодействие процессов (IPC)}
\label{ch:18-ipc}

\section{Зачем процессам общаться}

Каждый процесс изолирован: своя память, свои файловые
дескрипторы. Но реальные системы --- это множество
процессов, которым нужно обмениваться данными. Способы
такого обмена называют \textbf{IPC} (Inter-Process
Communication).

Все сценарии IPC выполняются на вашей ВМ на \fpath{pglab.nsu.ru}:
общаются процессы одной вашей системы.

\section{Файлы: самый простой способ}

Процесс пишет в файл, другой читает. Медленно, но
просто и наглядно. Для медленных потоков данных
(конфигурации, отчёты) этого достаточно. Ограничение:
файловая система не «сигнализирует» о новых данных ---
второму процессу нужно опрашивать.

\section{Пайпы (pipes)}

\textbf{Пайп} --- односторонний канал.
\textbf{Анонимный пайп} живёт в памяти и существует
только между связанными процессами (родитель и
потомком):
\begin{minted}[bgcolor=gray!6]{bash}
# оболочка построит этот конвейер за вас:
ls -l | grep txt | wc -l
\end{minted}

\textbf{Именованный пайп} (FIFO) --- файл-устройство
на диске, доступное любому процессу:
\begin{minted}[bgcolor=gray!6]{bash}
mkfifo myfifo
# в одном терминале:
cat > myfifo
# в другом терминале:
echo hello > myfifo
\end{minted}

Пайпы работают как буфер фиксированного размера
(обычно 64 КБ): если буфер полон, пишущий блокируется.

\section{Сокеты}

\textbf{Сокет} --- обобщение пайпа: двунаправленный
канал между двумя точками. Бывают:
\begin{itemize}
  \item \textbf{Unix domain} (межпроцессные, в пределах
    одной системы) --- быстрые;
  \item \textbf{TCP/IP} (между машинами по сети) ---
    так работают веб-сервер и клиент, SSH,
    большинство клиент-серверных приложений.
\end{itemize}

\section{Общая память (shared memory)}

Два процесса могут «пришить» одну и ту же физическую
страницу к своим адресным пространствам
(\nameref{ch:16-virtual-memory}) и обмениваться
данными без копирования. Самый быстрый способ IPC, но
требует синхронизации (семафоры, мьютексы), чтобы
процессы не записывали одновременно.

\section{Очереди сообщений}

Готовый механизм ядра: процесс кладёт сообщение
в очередь, другой --- забирает. Проще, чем общая
память (ядро следит за целостностью), медленнее.

\section{Сигналы}

Особый, асинхронный способ (\nameref{ch:12-processes}):
короткое уведомление без передачи данных
(«остановись», «перезапусти»). Сигналы --- это
уведомление, а не канал данных.

\section{Выбор способа}

\begin{center}
\begin{tabular}{p{0.52\linewidth}l}
\toprule
\textbf{Ситуация} & \textbf{Способ} \\
\midrule
одноразовые данные & файл \\
быстрый поток между родителем и дочерним & анонимный пайп \\
связь между несвязанными процессами & именованный пайп или сокет \\
клиент-сервер в одной системе & Unix socket \\
клиент-сервер по сети & TCP/IP socket \\
максимальная скорость & shared memory + синхронизация \\
\bottomrule
\end{tabular}
\end{center}

\section{Связь с остальным курсом}

Анонимные пайпы собирает оболочка
(\nameref{ch:07-text-pipelines}); сокеты лежат в основе
SSH (\nameref{ch:02-linux-access}) и Docker API
(\nameref{ch:19-docker-basics}); ограничения на
ресурсы контейнеров --- cgroups
(\nameref{ch:23-docker-resources}).

\section{Задачи для самопроверки}
\begin{enumerate}
  \item Перечислите основные способы IPC и
    краткое назначение каждого.
  \item Чем анонимный пайп отличается от именованного?
  \item Когда общая память лучше сокета, а когда ---
    нет?
  \item Какие данные передаёт сигнал?
  \item Почему shared memory требует синхронизации?
\end{enumerate}
